% ============================================================
%  chapters/ch06_conclusion.tex — Chapter 6: Conclusion
% ============================================================

\chapter{Conclusion}
% ------------------------------------------------------------
\section{Summary of Contributions}
Establishes an upstream computational audit methodology bridging data science and educational theory as a prerequisite benchmark for LLM deployment.

\section{Theoretical Implications}
Demonstrates that Multiplicity and Horizontality can be operationalized as scalable computational metrics in high-dimensional vector spaces.

\section{Future Research}
\subsection{Dimension B: Downstream Student Cognitive Agency}
Future studies will evaluate human-in-the-loop pedagogical dynamics and student reasoning.

\subsection{Multi-Turn Conversational Dynamics}
Extending audits to multi-turn interactions to analyze context drift over extended sessions.

\subsection{Expanded Model and Task Coverage}
Testing open-weight models and expanding task categories across multilingual contexts.

\section{Policy and Governance Implications}
Recommends mandatory institutional upstream audits and task-dependent deployment guidelines before integrating LLMs into active learning environments.
